\chapter{眼底中心凹定位}
\label{chap:fovea}

\section{任务、数据与度量}
本任务从彩色眼底图预测中心凹的横纵坐标。中心凹是黄斑区域的关键解剖位置，
可能缺少可直接阈值化的局部边界。高光、视网膜病变、视盘位置变化以及不同
成像尺度都会影响定位。课程数据包括 80 张带中心凹坐标的训练图和 20 张测试图，
属于 GAMMA 所涉及的图像域\citep{wu2023gamma}。原始 GAMMA 的任务范围还包括
青光眼分级和视盘视杯分割。本文只评价课程提供的中心凹坐标。
代码把标注访问限定在训练编号 1 至 80，课程测试编号 81 至 100 的公开同源
标签未纳入训练、校准和误差图。程序按内容哈希将完全相同的训练图绑定到同一折，
以免相同图像同时进入该折的训练与验证。

每个测试样本按 \code{ImageID,value} 分别写出 X、Y 坐标，共 40 行。
设真实和预测坐标分别为 $(x_i,y_i)$ 与 $(\hat x_i,\hat y_i)$，评价将
两轴元素合并计算
\begin{equation}
 \mathrm{RMSE}_{xy}
 =\sqrt{\frac{1}{2n}\sum_{i=1}^n
         \big[(\hat x_i-x_i)^2+(\hat y_i-y_i)^2\big]} .
 \label{eq:fovea-rmse}
\end{equation}
单位为原始图像像素。单张图的欧氏距离除以 $\sqrt 2$，才得到与该公式一致的
单图坐标 RMSE。公开榜和私有榜分别使用约一半测试图，任何少量大误差都会影响
平方和。本文在同一张表中呈现两侧结果并保留各自的样本口径。

\section{热图建模与坐标恢复}
\subsection{中心裁剪与血管图}
程序将眼底原图做中心方形裁剪，避免将非方形眼底直接挤压到方形。
设裁剪区域左上角为 $(x_0,y_0)$、边长为 $S$，则归一化坐标为
$u=(x-x_0)/S$、$v=(y-y_0)/S$。预测后用相反变换恢复原图坐标。
几何增强使用同一坐标变换处理图像、血管图与标注点。裁剪可能移除边缘结构，
这是在保持纵横比例和保留完整视野之间的取舍。

在 RGB 之外增加由 Fundus Image Toolbox 生成的血管图，形成四通道输入。
视网膜血管弓的走向与中心凹的空间关系提供额外位置线索，即使单个像素
的颜色不足以指示凹陷区域。训练与推理均对血管图使用对应的中心裁剪和尺寸。
该输入来自外部血管模型的图像推理，血管响应在病变或弱对比图像中也会有误差，
因此保留 RGB 主输入。

\archdiagram{fovea_arch}{中心凹热图网络与坐标级融合}{fig:fovea-arch}

\subsection{编码器、解码器与定位目标}
主线采用 EfficientNet-B1，输出步幅为 4、8、16、32 的四级特征。U 形解码器
逐级对深层特征做双线性插值，与同尺度编码器特征拼接，三个解码块的输出通道数
依次为 128、64、32。每个块包含两组卷积、批归一化与 SiLU。
$1\times1$ 卷积和 Sigmoid 产生 $128\times128$ 单通道热图。相比
全局坐标回归，空间监督可保留多个候选区域和中心凹的相对位置。视盘辅助头、
坐标回归头和 VascX 第五输入通道均是候选路径，主线配置关闭这些选项。

标注坐标在热图格点上形成带软边缘的圆盘，半径为 6 格。对距离中心为 $d$ 的
位置，代码目标近似为 $g=\sigma((6-d)/0.75)$。虽然运行配置保留通用的
\code{loss=mse} 字段，实际代码在 \code{target=circle} 时优先进入
BCE、Dice 与坐标正则分支。令 $p_j$ 和 $g_j$ 分别为预测与目标热图像素，
训练目标为
\begin{equation}
 \mathcal L =
  \mathrm{BCE}(\mathbf p,\mathbf g)+
  \left(1-\frac{2\sum_j p_jg_j+1}{\sum_j p_j+\sum_j g_j+1}\right)
  +0.05\,\big\|\operatorname{centroid}_{p^2}(\mathbf p)-(u,v)\big\|_1 ,
 \label{eq:fovea-loss}
\end{equation}
其中 Dice 项和坐标误差在批内求均值，坐标 L1 在两轴上取均值。
平方权重使训练用质心更关注高响应区域。推理时先找热图最大响应位置，再在
以峰值为中心的 17 格窗口内计算归一化质心，降低远端次峰的影响。镜像预测在
统一方向下合并后恢复至原图坐标。

\section{训练过程与融合选择}
主线的种子 123 与 777 各训练一套 10 折模型。训练端使用 AdamW、
AMP 和梯度裁剪，配置为批大小 8、60 轮、学习率 $10^{-4}$、
权重衰减 $10^{-4}$、3 轮预热和后续余弦调度。各折根据验证坐标
RMSE 保存最优模型。种子 123 的配置文件记录了补跑折的范围，
全部十折是否完成应以各自独立的日志为准。训练与验证的误差指标不同。
训练优化式~\eqref{eq:fovea-loss}，验证使用恢复到原图后的
式~\eqref{eq:fovea-rmse}。

\chart{fovea_training.pdf}{种子 123 第 0 折的热图训练与原图像素验证}{fig:fovea-curve}
{RMSE 用对数纵轴展示训练早期的大幅变化。标记点是日志中 RMSE 最低的轮次。}

在单独评估 VascX 分支时，项目使用中心方形裁剪、1024 输入及 512 滑窗。
VascX 论文描述眼底血管与关键点模型，本文的本地推理适配和融合规则依赖
项目记录\citep{vargas2025vascx}。先合并种子 123 和 777 的本地模型预测，
再根据本地坐标和 VascX 坐标的距离分档融合。设本地预测为 $\mathbf a$、
VascX 预测为 $\mathbf b$，融合可写为
\[
 \hat{\mathbf x}=\bigl(1-w(d)\bigr)\mathbf a+w(d)\mathbf b,\qquad
 d=\|\mathbf a-\mathbf b\|_2 .
\]
早期 v12 在允许融合时使用 20\% 的外部模型权重。v15 使用多个距离区间，
历史记录保留了评价与提交文件，具体分档阈值及完整生成命令未在当前代码
目录中找到。附录提供两条坐标生成入口，最终提交的分数用 Kaggle 记录核对。

\section{结果、负面结果与典型误差}
\subsection{本地评估与公开、私有榜}
表~\ref{tab:fovea-local} 分列单种子、固定权重和距离分档结果。
种子 123 和 777 的固定融合在本地达到 13.357 px，
使用局部共识保护加入 VascX 后为 13.197 px。距离分档在拟合使用的
样本上得到 12.994 px，嵌套验证为 13.190 px。后者与 v12 的差距很小，
支持谨慎解释该增益。种子 2026 的独立模型组在本地为 15.022 px，
加入固定权重的搜索未改善主线。

\begin{table}[H]
\centering\small
\caption{中心凹定位本地验证与融合评估}\label{tab:fovea-local}
\begin{tabularx}{\linewidth}{@{}Ylr@{}}
\toprule
方案 & 评价方式 & RMSE（px） \\
\midrule
种子 123，十折血管输入 & OOF & 13.593 \\
种子 777，十折血管输入 & OOF & 14.174 \\
种子 2026，十折血管输入 & OOF & 15.022 \\
v11，种子 123/777 固定融合 & OOF & 13.357 \\
v12，固定融合加 VascX & OOF & 13.197 \\
v15，距离分档 & 样本内 & 12.994 \\
v15，距离分档 & 嵌套验证 & 13.190 \\
\bottomrule
\end{tabularx}
\end{table}

\begin{table}[H]
\centering\small
\caption{中心凹定位的 Kaggle 提交，RMSE 越低越好}\label{tab:fovea-score}
\begin{tabularx}{\linewidth}{@{}Yrr@{}}
\toprule
提交 & Public（px） & Private（px） \\
\midrule
v2，B0 与 B1 基线集成 & 13.65941 & 25.63019 \\
v7，较早的血管输入方案 & 13.29574 & 22.75990 \\
v11，两种子固定融合 & 11.64967 & 21.72259 \\
v12，局部共识融合 & 10.35936 & 20.16412 \\
v15，距离分档融合 & \textbf{9.90927} & \textbf{19.54239} \\
\bottomrule
\end{tabularx}
\end{table}

v15 相对 v2 的公开榜下降 3.75014 px，私有榜下降 6.08780 px。
同一 v15 在私有榜比公开榜高 9.63312 px，说明两侧样本误差分布存在
明显差异。表~\ref{tab:fovea-score} 给出历史提交间的观察对比，
图~\ref{fig:fovea-results} 展示分折误差和两侧榜单记录。
两侧分数在这些阶段呈同方向下降，差值涉及血管输入、折数和融合规则，
不能拆解为某个模块的独立效果。

\chart{fovea_results.pdf}{种子 123 的逐折误差与提交演变}{fig:fovea-results}
{左图用每折保存的 OOF 坐标重新计算，分折样本数为 7 至 9 张。右图使用 Kaggle 返回的两侧榜单成绩。}

\clearpage
\subsection{误差集中与未采用方案}
种子 123 的第 7 折 RMSE 约 30.99 px，其中单个样本占该折坐标平方误差
约 80.2\%。第 8 折约为 5.06 px。对少样本任务，报告总体 RMSE 的同时
给出分折值与误差集中度，便于判断融合对少数极端样本的影响。
\ifLocalReport 图~\ref{fig:fovea-case} 使用验证折保存的坐标说明成功与偏移情况。
中心裁剪中的坐标还原、热图次峰与血管预测偏差均可造成偏移，
仅凭单张图无法分离这些因素。\fi

\privatechart{fovea_case.png}{中心凹验证案例在原图坐标中的对照}{fig:fovea-case}
{绿色圆标记人工中心凹，橙色十字标记折外预测。每幅图的误差按两轴平方均值开根计算。}

REFUGE 领域预训练、VascX 先验作为第五输入通道、困难样本加权及裁剪
平移 TTA 都在部分固定折验证中尝试过。记录中的困难折结果显示，额外
先验输入的误差约 34.9 px，明显弱于同折血管输入方案。全局融合 VascX
的收益有限，将它保留为独立坐标分支，使融合可利用其强项并减弱高分歧时的影响。
这些试验基于有限折，足以解释未纳入主方案的原因。

\section{本章小结}
本任务将像素级热图与解剖血管信息结合，通过局部质心恢复原图坐标。
按预测距离使用外部分支后，v15 在已核验提交中同时获得最低公开与私有
RMSE。较大的榜单分差和单样本误差占比说明，独立样本上的定位稳定性仍是
主要局限。
